\input{preamble}
\input{glossary}
\addbibresource{bibfile.bib}


\usepackage[utf8]{inputenc}
\usepackage{mathrsfs}
\newtheorem{thm}{Théorème}
\newtheorem*{thm*}{Théorème}

% remarks in red
\usepackage{xcolor}
\newcommand{\rmq}[1]{\textcolor{red}{#1}}
\newcommand\ignore[1]\null

% \setcounter{secnumdepth}{0} % sections will not be enumerated

\usepackage[figurename=Fig.]{caption}
\usepackage[tablename=Tab.]{caption}


\usepackage{fancyhdr,graphicx,lastpage}% http://ctan.org/pkg/{fancyhdr,graphicx,lastpage}
\fancypagestyle{plain}{
  \fancyhf{}% Clear header/footer
  \fancyhead[R]{\textit{Internship report}}% Right header
  \fancyhead[L]{\textit{Your name}}% Right header
  \fancyfoot[L]{\textit{University year 20xx-20xx}}% Left footer
  \fancyfoot[R]{\thepage\  / \pageref{LastPage}}% Right footer
}
\pagestyle{plain}% Set page style to plain.



\begin{document}
\begin{sloppypar}
\pagestyle{empty}

% Declare new goemetry for the title page only.
\newgeometry{margin=1in}
\newcommand\titleofdoc{\bfseries Internship report M1 (2A ENSIIE)}

 %%%%% Put your document title in this argument
\newcommand\GroupName{Group Name} %%%%% Put your group name here. If you are the only member of the group, just put your name
\begin{titlepage}
   \begin{center}
        \vspace*{2cm} % Adjust spacings to ensure the title page is generally filled with text
        %% input logos
        {\includegraphics[width=0.22\textwidth]{images/logo/univ-rouen.jpg} \quad \quad}{\includegraphics[width=0.2\textwidth]{images/logo/lmrs.png} \quad \quad}
        {\includegraphics[width=0.11\textwidth]{images/logo/cnrs.png}}\\

        \noindent\rule[0.25\baselineskip]{\textwidth}{1pt}
        
        \vspace{0.5 cm} 

        \Huge{\titleofdoc} 

        \vspace{0.5cm}
        \LARGE{Title / subject of the internship}
          
        \vspace{2 cm}
        % \Large{\GroupName}
        \vspace{0.25cm}
        \large{Your name here}\\
        \large{\bfseries Internship carried out from xx/xx/20xx  to  xx/xx/20xx}\\
        % \vfill
        \vspace{4 cm}
         
        \leftline{ \small { \textbf{Internship tutors :} M. xxx}}
        \leftline{ \small { \textbf{Referent teacher :} M. xxx}}
        \vspace{1 cm}
        \leftline{\small { \textbf{\'Educational institution :} Université Paris-Saclay / ENSIIE - M1 Maths Appliquées}}
        \leftline{\small { \textbf{Internship host company :} LMRS - Av. de l'Université, BP.12, 76801 Saint-Étienne-du-Rouvray}}
        
        \noindent\rule[0.25\baselineskip]{\textwidth}{1pt}\\
        
        {\includegraphics[width=0.11\textwidth]{images/logo/ensiie.png}\quad}
        % {\includegraphics[width=0.4\textwidth]{images/logo/datalab.jpeg} \quad}
        {\includegraphics[width=0.23\textwidth]{images/logo/saclay.png}\quad}\\
      
    \end{center}
\setcounter{page}{0}
\end{titlepage}
%%% notice: 
% . prénom et nom de l’élève
% . titre du stage
% . type du stage (1A, 2A ou 3A + humanitaire ou autre mise en situation
% pro. . .)
% . nom de l’établissement
% . prénom et nom de l’encadrant·e dans l’établissement
% . dates du stage


%%% Ends the declared geometry for the titlepage
\restoregeometry


\centerline{\huge {\textbf{Acknowledgments} } }
\setcounter{page}{-5}

% Acknowledgments
Say something here.


% \newpage 
% \ % The empty page
\newpage





% \section{Préambule}
\begin{abstract}
\noindent
Write an abstract here. \\  \\
\textit{Key words}: financial models, machine learning...
\end{abstract}








\newpage
\tableofcontents




% Corps du rapport
\newpage
\setcounter{page}{1}
\pagestyle{plain}
\section{Introduction} 
%%% use glossary like this: 
\Gls{put} and \gls{call} options are important in finance.

% (0.5 - 1 page maximum, accessible to all)
% • Description ciblée de l’environnement
% • Description ciblée du sujet et des objectifs du stage
% • Description ciblée des moyens mis en œuvre
% • Présentation contextualisée du plan du rappor


\section{Work environment}
\subsection{The host company}
% description de l’établissement, du service, de l’équipe, du client


\subsubsection{Different functionality}
\paragraph{You can add a figure like this:} 
%%% add figure: 
\begin{figure}[H] %% H: fixed; h: flexible
  \centering
  \includegraphics[scale = 0.2]{images/logo/ensiie.png}
  \caption{(Put a figure description here.)}  
  \label{fig:ensiie} %% figure reference name, can be cited after. 
\end{figure}

\paragraph{You can also add an equation: }
%%% add an equation: 
\begin{equation} 
     dS_t=\mu S_t\text{d}t+\sigma S_t\text{d}B_t ,
\label{eq:bs-dynamique} %% equation reference name, can be cited after. 
\end{equation} 

%% explain iems: 
and explain it: 
\begin{itemize}[label=$\bullet$]
\setlength\itemsep{-0.5em}
\item $t$: temps, en années ; 
\item $S_t$: prix de  l’actif sous-jacent $S$ à l'instant $t$ ; 
\item $B_t$: un mouvement brownien à l'instant $t$ ; 
\item $\mu$ : drift de $S$ ; 
\item $\sigma$ : la volatilité (constante) du sous-jacent.
% \item $r$: le taux d'intérêt sans risque  d'un emprunteur sûr  dont la rentabilité est certaine sur une période donnée.
\end{itemize}

\paragraph{Moreover, you can add a table: }
%% add a table
\begin{table}[H]
    \centering
    \begin{tabular}{lrr}
    \hline
        Méthode & Temps de calcul (en s) & Prix de call \\
    \hline
        line 1 & 1 & 2 \\
        line 2 & 1 & 2 \\
        line 3 & 1 & 3 \\
        line 4 & 1 & 9 \\
    \hline
    \end{tabular}
    \caption{Performance des méthodes de pricing du modèle de Heston}
    \label{tab:heston_methods} %% table reference name, can be cited after. 
\end{table}

\paragraph{These things can be cited as below: }

%%% citation example: 
The Fig.\ref{fig:ensiie} used the equation (\ref{eq:bs-dynamique}), and turned out the results as shown in Tab.\ref{tab:heston_methods}.




\subsection{The host team}
















\newpage
\section{(General problematic statement)}
%%% notice: 
% . description du contexte et des objectifs
% . Quel est le travail à réaliser ? À qui est-il destiné ?
% . mise en perspective des enjeux du stage
% Pourquoi est-ce important de réaliser ce travail ?
% . description des étapes nécessaires pour atteindre les objectifs
% Comment s’articulent ces différentes étapes ? Qui va les réaliser ?
% . description de vos mission




\newpage
\section{(Your contribution)}
\subsection{State of the art}

\subsection{Tools or conception used}

\subsection{Realisation}

\subsection{Results and discussion}









\newpage
\section{Conclusion et perspectives}
\label{conclusion}
\cite{perceptron}




\newpage

\printbibliography % [title={Whole bibliography}]


% \bibliographystyle{IEEEtran}
% \bibliography{bibfile}




\newpage
\label{glossaire}
%  vocabulaire technique, avec définition et page où est utilisé le mot la première fois

\setglossarystyle{listhypergroup}
% \setglossarystyle{index} % not combine by group
\glsaddallunused % change line
\printnoidxglossaries





\newpage
\appendix % macro

\section{Annexe DD\&RS}
\label{ddrs}
%% write as you write sections

\subsection{Sustainable Development }


\subsection{Social Responsibility }


% \newpage
% \section{Codes}


% % add codes
% \subsection{Traitement et Audit des Données par séries temporelles}
% \inputminted[linenos=true, breaklines=true]{python}{codes/time_series_analysis.py}

\end{sloppypar}

\end{document}
